\subsection{Field generation / reconstruction packages (T3-2)}
\label{app:field_reconstruction_theorems}

\AuditTag This appendix records the strong-closure bridge target ``net$\to$field'' (and, when needed, ``field$\to$net'') in a theorem-facing way.
In 4D interacting settings, reconstruction typically requires additional regularity assumptions; the manuscript therefore packages them explicitly.

\subsubsection{Condition bundle: reconstruction regularity}
\begin{assumption}[Reconstruction condition bundle (template)]
\label{ass:reconstruction_bundle_template}
Assume the following for a local net $(\mathcal{A}(\mathcal{O}))$ in a vacuum representation:
\begin{itemize}
\item \textbf{(RC1) Domain control:} Assumption~\ref{ass:common_invariant_domain_generators} holds for the intended generators.
\item \textbf{(RC2) Localization regularity:} there exists a generating family compatible with the region localization (pointlike or stringlike, as declared).
\item \textbf{(RC3) Energy bounds:} generator norms on energy-bounded vectors satisfy polynomial energy bounds ensuring temperedness of correlations.
\end{itemize}
\end{assumption}

\subsubsection{Reconstruction theorem (template)}
\begin{theorem}[Net$\to$field reconstruction (template)]
\label{thm:net_to_field_reconstruction_template}
\MathTag Under Assumption~\ref{ass:reconstruction_bundle_template}, there exists a family of (operator-valued distribution) fields $\Phi$ whose bounded functional calculus generates the net, and whose vacuum correlation functions satisfy the Wightman structural package at the declared regularity level.
\end{theorem}

\AuditTag Theorem~\ref{thm:net_to_field_reconstruction_template} is a theorem-facing target statement.
If any of (RC1)--(RC3) cannot be discharged, then the failure is recorded by the canonical W1/W2 records (Appendix~\ref{app:minimal_failure_point_templates}), and the manuscript does not claim field reconstruction.

